\newpage
\phantomsection
\label{prf:euclid-i-2}
\LRAProofFor{thm:euclid-i-2}

\begin{remark*}[Return]
\hyperref[thm:euclid-i-2]{Return to Theorem}
\end{remark*}

\begin{theorem*}[Euclid I.2: Copying a segment from a given point]
Given a point $A$ and a finite straight line segment $BC$, there exists a
straight line segment beginning at $A$ and equal to $BC$.
\end{theorem*}

\begin{proof}[Professional Standard Proof]
\LRAProofBodyStart
Let $A$ be the given point and let $BC$ be the given segment. Join $AB$ and
construct an equilateral triangle $DAB$ on $AB$ by Euclid I.1. Produce the
straight lines $DA$ and $DB$. With center $B$ and distance $BC$, describe a
circle meeting the produced line $DB$ at $G$. With center $D$ and distance
$DG$, describe a circle meeting the produced line $DA$ at $L$.

Since $BG=BC$ and $DL=DG$, while $DA=DB$ because $DAB$ is equilateral, the
remainders $AL$ and $BG$ are equal. Therefore $AL=BC$. Thus from the point $A$
a straight line segment equal to the given segment $BC$ has been constructed.
\end{proof}

\begin{proof}[Detailed Learning Proof]
\LRAProofBodyStart
\begin{center}
\begin{tikzpicture}[scale=1.15, line cap=round, line join=round]
  \definecolor{euclidblue}{RGB}{30,110,190}
  \tikzset{
    equaltick/.style={
      postaction={decorate},
      decoration={markings, mark=at position 0.5 with {
        \draw[line width=0.6pt] (0,-2.4pt)--(0,2.4pt);
      }}
    }
  }

  \coordinate (A) at (0,0);
  \coordinate (B) at (2.1,0);
  \coordinate (C) at (3.0,0);
  \coordinate (D) at (1.05,1.65);
  \pgfmathsetmacro{\seg}{0.9}
  \pgfmathsetmacro{\db}{sqrt((2.1-1.05)^2+(0-1.65)^2)}
  \pgfmathsetmacro{\scaleG}{\seg/\db}
  \coordinate (G) at ($(B)+\scaleG*(B)-\scaleG*(D)$);
  \pgfmathsetmacro{\dg}{\db+\seg}
  \coordinate (L) at ($(D)!\dg/\db!(A)$);

  \draw[gray!50, thin] (B) circle (\seg);
  \draw[gray!45, thin] (D) circle (\dg);

  \draw[gray!70, thin] (D) -- ($(D)!1.45!(A)$);
  \draw[gray!70, thin] (D) -- ($(D)!1.45!(B)$);
  \draw[gray!70, thin] (A) -- (B);
  \draw[gray!70, thin] (D) -- (A) -- (B) -- cycle;

  \draw[very thick, equaltick] (B) -- (C);
  \draw[very thick, euclidblue, equaltick] (A) -- (L);
  \draw[very thick, euclidblue] (B) -- (G);
  \draw[very thick, euclidblue] (D) -- (G);
  \draw[very thick, euclidblue] (D) -- (L);

  \foreach \p/\pos in {A/below,B/below,C/below,G/below right,D/above,L/below left}{
    \fill (\p) circle (0.023);
    \node[\pos] at (\p) {$\p$};
  }

  \node[font=\footnotesize] at ($(B)!0.5!(C)+(0,-0.28)$) {$BC$};
  \node[euclidblue, font=\footnotesize] at ($(A)!0.5!(L)+(-0.20,-0.16)$) {$AL=BC$};
  \node[gray!70, font=\footnotesize] at ($(B)+(0.65,-0.70)$) {center $B$};
  \node[gray!70, font=\footnotesize] at ($(D)+(-1.00,1.05)$) {center $D$};
\end{tikzpicture}
\end{center}


The role of the equilateral triangle is to create two matching initial lengths:
$DA=DB$. After the side $DB$ is extended, the circle centered at $B$ marks a
point $G$ with $BG=BC$. Then the circle centered at $D$ transfers the full
length $DG$ onto the other produced side, giving $DL=DG$.

Now compare the two long segments from $D$. On one side,
$DG=DB+BG$; on the other, $DL=DA+AL$. Since $DG=DL$ and $DB=DA$, subtracting
the equal initial parts leaves $BG=AL$. But $BG=BC$, so $AL=BC$ by Common
Notion I. The segment $AL$ begins at the required point $A$ and has the
required length.
\end{proof}

\begin{remark*}[Proof structure]
Euclid I.2 is a transfer construction. It does not assume a compass can be
lifted and moved directly; instead it copies the length $BC$ using only
straight lines, produced straight lines, circles, Euclid I.1, and common
notions about equality.
\end{remark*}

\begin{dependencies}
\begin{itemize}
  \item \hyperref[def:euclid-point]{Point}
  \item \hyperref[def:euclid-straight-line]{Straight line}
  \item \hyperref[def:euclid-circle]{Circle}
  \item \hyperref[thm:euclid-i-1]{Euclid I.1}
  \item \hyperref[ax:euclid-postulate-draw-straight-line]{Euclid Postulate I}
  \item \hyperref[ax:euclid-postulate-produce-straight-line]{Euclid Postulate II}
  \item \hyperref[ax:euclid-postulate-describe-circle]{Euclid Postulate III}
  \item \hyperref[ax:euclid-common-notion-things-equal-to-same]{Euclid Common Notion I}
\end{itemize}
\end{dependencies}

\clearpage
